% !TeX root = ../main.tex
\section{Discussion}
\label{sec:discussion}

\subsection{Principal Accuracy Bottlenecks}

The combined ResNet, ConvNeXt, and Tiny-ViT analyses indicate that the dominant limitations under the canonical deployment regime (standard 4-bit quantization with 5\% C2C and 10\% D2D variability) are not the ones most often emphasized at the device level. Once hardware-aware training (HAT) and scale recovery are available, standard cycle-to-cycle (C2C) and device-to-device (D2D) variability do not dominate the error budget. Three constraints dominate instead.

A quantitative decomposition supports this ordering. First-order Sobol indices computed from a $7 \times 9$ D2D--ADC grid show that, over the full parameter space, ADC precision accounts for 97.6\% of accuracy variance ($S_{\text{ADC}}=0.98$), while D2D contributes only 1.8\%. However, conditioning on the operational envelope ($\geq$6-bit ADC, $\sigma_{\text{D2D}} \leq 15\%$), the hierarchy inverts: D2D variability accounts for 92.2\% of residual variance ($S_{\text{D2D}}=0.92$). This two-phase structure has a direct hardware implication: designers should first secure 6-bit readout, then invest the remaining error budget in reducing device-to-device mismatch.

The first is ADC resolution. The transition around 6-bit ADC resolution marks a clear change in model behavior: below this point transformer inference degrades substantially, indicating that improvements in conductance control remain gated by readout precision in the present mixed-signal stack.

The second is hardware-instance overfitting. The collapse of V4 under fresh-instance transfer (10.00\% accuracy) shows that standard HAT readily fits a single fixed D2D realization. Resampling D2D masks during training raises fresh-instance accuracy to \(86.37 \pm 1.54\%\), indicating that multi-instance-aware training can mitigate this transfer failure mode. Mechanistically, Ensemble HAT changes the fixed spatial mismatch map itself across epochs rather than merely perturbing one nominal instance.

A residual gap to the 98.06\% digital baseline persists under matched deployment conditions. Part of this difference may reflect regularization from per-epoch D2D resampling rather than analog noise alone. Capacity-aware Ensemble HAT schedules remain an open optimization problem.

The third is the scale-masking regime observed in V2. The high accuracy of V2 under standard uniform noise is a conditional property of the canonical simulator configuration rather than a universal form of robustness. When digital scale recovery maps analog perturbations below the 4-bit quantization step, many realizations remain in the same recovered weight bin after calibration. This protection depends on accurate post-array gain calibration and disappears once the noise law becomes state-dependent, as shown by the proportional-noise experiments.

\subsection{Transformer Sensitivity to Non-Idealities}

Comparing the dual testbeds—ConvNeXt (CNN) and Tiny-ViT (Transformer)—reveals that the Transformer architecture is more fragile under the present physical stress tests. Because the two models serve complementary roles (ConvNeXt from scratch versus Tiny-ViT from pre-trained fine-tuning), this result should be read as evidence of a strong architecture-and-training interaction rather than as a universal ranking of all Transformers against all CNNs.

The proportional-noise experiments illustrate this most clearly. The collapse of V4 under proportional noise (10.00\% accuracy), contrasted with the relative stability of ConvNeXt under moderate non-linearities, indicates that Transformer inference depends on the precision of token projections. Because attention maps arise from global interactions, perturbations at a small number of high-conductance devices can disrupt the spatial weighting of the entire block. The optical front-end experiments point in the same direction: the larger accuracy drop in V6 than in ResNet R4 implies that input distortions are propagated more readily through sequence-wide attention than through a purely convolutional hierarchy.

Robustness is highly regime-matched. Tiny-ViT can recover strongly when trained and evaluated under the same proportional-noise law, yet the resulting checkpoint does not transfer back to the canonical uniform-noise regime. In this framework, robustness is distribution-specific rather than noise-invariant.

The framework remains reasonably resilient overall. A compound stress test combining all modeled non-idealities (2\% C2C, 3\% D2D, 6-bit ADC, 1\% IR drop, 1\% sneak path, and $1{,}000$ s retention) maintained an 89.61\% accuracy for Tiny-ViT, suggesting that the mixed-precision architecture with Ensemble HAT can tolerate moderate cumulative deployment stress within the present model.

\subsection{Task Complexity and Data Starvation}

The multi-dataset results show a clear interaction between task difficulty and analog robustness. HAT recovery is most valuable on CIFAR-100, where the decision boundaries are comparatively fine, whereas CIFAR-10 is sufficiently simple that even noisy standard models remain functional. However, ResNet-18 exhibits a catastrophic failure mode on CIFAR-100 under the canonical noisy regime, collapsing to chance-level accuracy (1.00\%) regardless of HAT training. A post hoc audit showed that the earlier CIFAR-10 inconsistency arose from a legacy model-loading issue rather than from the analog inference path itself. Replay with the compatibility loader recovered the two affected CIFAR-10 checkpoints to 94.12\% and 89.60\%. Table~\ref{tab:result-summary} retains the original checkpoint-best convention for the canonical R4 row (90.37\%), whereas the discussion quotes the post-fix replay value because that audit established that the saved weights were valid and that the earlier collapse was a load-time artifact. The CIFAR-100 result, by contrast, remained at 1.00\% after this correction and was therefore confirmed as a genuine optimization failure rather than a load-time artifact. Because this failure reflects a limitation of the present ResNet conversion and training recipe rather than a robust cross-architecture physical trend, we do not draw architectural conclusions from the ResNet-18 CIFAR-100 result and focus the cross-architecture comparison on Tiny-ViT and ConvNeXt, which both pass baseline sanity checks under the same workflow. Flowers-102 marks the opposite extreme for the current recipe: Tiny-ViT recovers only partially, and the ConvNeXt control remains unstable in this low-data regime. Because the ConvNeXt Flowers-102 baseline is only a single-run boundary estimate, we treat that dataset as a qualitative low-data stress test rather than as a stable cross-architecture ranking benchmark.

\subsection{The "Analog Ceiling" and Energy Efficiency}

A first-order analytical energy model projects a potential 11.45× reduction in dense-projection energy relative to FP32 digital inference, though this illustrative upper bound rests on unvalidated edge-node placeholders. Moderate routing overhead (10–50\%) reduces this gain to 9.90–11.10×, and digital execution of dynamic attention still dominates the overall footprint (57.9\% of total energy). Even so, the estimated 273.94~\(\mu\)J footprint of the hybrid stack compares favorably with reported INT8 ViT digital accelerators \citep{li2023ivit}, suggesting that hybrid deployment may offer an efficiency advantage in favorable regimes. These numbers should be interpreted as first-order system-level upper bounds under the current placeholder constants rather than as chip-predictive routed-silicon estimates.

\subsection{Limitations}
\label{subsec:limitations}

Several physical idealizations of the behavioral model should be noted. The energy model parameters ($E_{cell}$, $E_{ADC}$, $E_{DAC}$) are analytical placeholders rather than chip-predictive measured values. Adding extra routing overhead equal to 10\%, 30\%, or 50\% of the analog-MAC budget would reduce the referenced gain, indicating that the qualitative advantage is robust but sensitive to interconnect assumptions.

Regarding array non-idealities, our sensitivity sweep for position-dependent IR drop (up to 3\%) and sneak-path currents (up to 2\%) suggests robustness to moderate parasitic effects. However, these baseline ranges are derived from metallic-interconnect ReRAM literature \citep{fastirdrop2025,iconniv2025}. Given that organic semiconductor arrays often exhibit significantly higher sheet resistances, this range should be interpreted as a lower-bound sensitivity probe rather than a prediction of full-scale array behavior. Future integration of measured organic-array sheet-resistance data will be necessary to refine these bounds. At the same time, the accompanying supplementary proxy-sensitivity sweeps and parameter-risk tables indicate that, within the tested uncertainty ranges, proxy-value perturbations do not reverse the main ranking of ADC-limited versus D2D-limited regimes. Furthermore, the observed "scale-masking" effect, where small C2C noise is absorbed within quantization bins, relies heavily on the assumption of ideal calibrated digital rescaling. A full-test hook-based check reported in the Supplementary Information shows that moderate ADC calibration errors (\(\pm 0.5\) LSB offset, \(\pm 5\%\) gain error, \(0.5\) LSB INL) change the canonical Tiny-ViT V4 accuracy by only \(0.18\) percentage points, whereas a deliberately pessimistic combination causes a \(5.14\) percentage-point drop. In practical deployment, readout-chain non-idealities therefore remain a gate-keeping condition, but the protected regime is not immediately destroyed by moderate calibration error.

Finally, the observed degradation at $NL=2.0$ ($27.72\pm0.82\%$) should be interpreted as a recovery limit for the present gradient-scaling surrogate and training recipe, rather than as an immutable physical constraint of organic optoelectronic materials themselves. A supplementary group-wise gradient diagnostic on the frozen Tiny-ViT V4 checkpoint further localizes this surrogate failure: under matched forward conditions, enabling $NL=2.0$ only in the MLP analog blocks lowers the affected-parameter gradient cosine to $0.815$ and the norm ratio to $0.671$, whereas patch embedding and attention projections remain essentially unchanged. This indicates that the present nonlinear-write bottleneck is concentrated in the MLP path rather than being uniformly distributed across the transformer. Alternative straight-through formulations or specialized non-linear-aware optimizers may shift this boundary and remain an open research direction. Taken together, these simplifications mean that the framework should be read as a materials-to-system decision aid for ranking deployment risks under plausible parameter ranges, not as a chip-predictive emulator for a specific fabricated array.

\subsection{Outlook}

An immediate next step is to close the loop between laboratory characterization and simulation through automated profile fitting. Multi-instance HAT schedules should also be explored further, particularly where transferability and capacity appear to compete. A preliminary cross-framework comparison against CrossSim shows reasonably consistent baseline inference under clean conditions (86.2\% ours vs.\ 83.7\% CrossSim on a 1\,000-sample subset at 8-bit ADC), but the accuracy gap widens substantially under noise injection (81.6\% vs.\ 67.2\% at $\sigma=5\%$), highlighting the sensitivity of accuracy predictions to the specific noise-to-conductance mapping. Because the two toolchains use different internal approximations for mapping device assumptions to effective conductances, this discrepancy should be read as a joint-calibration problem rather than as evidence that one simulator is categorically correct and the other is not. These shared-regime sanity checks can be extended into broader numerical-equivalence studies that also include retention, photoresponse, and richer write dynamics.
